\section*{Capítulo \ref{ch:g10:lewis-and-shape} --- Las estructuras de Lewis y la geometría de las moléculas}

\begin{solution}{exo:g10:lewis-and-shape:1}
Un par de \omterm{def:g9:inside-the-atom:nucleus}{electrones} compartido por dos \omterm{def:g7:atoms-and-molecules:atom}{átomos}. Un par de \omterm{def:g10:electron-shells:valence}{electrones de
valencia} que pertenece a un solo \omterm{def:g7:atoms-and-molecules:atom}{átomo}, sin compartir.
\end{solution}

\begin{solution}{exo:g10:lewis-and-shape:2}
H: 1 (le falta un \omterm{def:g9:inside-the-atom:nucleus}{electrón} para su dúo). C: 4, N: 3, O: 2 (les faltan 4,
3 y 2 \omterm{def:g9:inside-the-atom:nucleus}{electrones} para el octeto).
\end{solution}

\begin{solution}{exo:g10:lewis-and-shape:3}
\chemfig{H-\charge{0=\:,90=\:,270=\:}{Cl}}: el cloro tiene el par compartido y tres \omterm{def:g10:lewis-and-shape:lone-pair}{pares
no enlazantes}, 8 \omterm{def:g9:inside-the-atom:nucleus}{electrones}.
\end{solution}

\begin{solution}{exo:g10:lewis-and-shape:4}
$5 + 3 \times 1 = 8$ \omterm{def:g9:inside-the-atom:nucleus}{electrones}, 4 pares (tres enlaces y un \omterm{def:g10:lewis-and-shape:lone-pair}{par no
enlazante}).
\end{solution}

\begin{solution}{exo:g10:lewis-and-shape:5}
\ce{CH4}: tetraédrica. \ce{NH3}: piramidal. \ce{H2O}: angular.
\ce{CO2}: lineal.
\end{solution}

\begin{solution}{exo:g10:lewis-and-shape:6}
\chemfig{H-\charge{90=\:,270=\:}{S}-H}: como el agua, cuatro grupos alrededor del azufre,
de los que dos son \omterm{def:g10:lewis-and-shape:lone-pair}{pares no enlazantes}: angular.
\end{solution}

\begin{solution}{exo:g10:lewis-and-shape:7}
$2 \times 5 = 10$ \omterm{def:g9:inside-the-atom:nucleus}{electrones}, 5 pares. Un enlace sencillo y 4 \omterm{def:g10:lewis-and-shape:lone-pair}{pares no
enlazantes} dejan a cada nitrógeno con 6 \omterm{def:g9:inside-the-atom:nucleus}{electrones}; convirtiendo dos
\omterm{def:g10:lewis-and-shape:lone-pair}{pares no enlazantes} en pares compartidos se obtiene un \omterm{def:g10:lewis-and-shape:multiple-bond}{enlace triple} y
un \omterm{def:g10:lewis-and-shape:lone-pair}{par no enlazante} en cada \omterm{def:g7:atoms-and-molecules:atom}{átomo}:
\chemfig{\charge{180=\:}{N}~\charge{0=\:}{N}}.
\end{solution}

\begin{solution}{exo:g10:lewis-and-shape:8}
\chemfig{H-C(-[2]H)(-[6]H)-\charge{90=\:,270=\:}{O}-H}: cuatro enlaces alrededor del
carbono, tetraédrica; alrededor del oxígeno, dos enlaces y dos \omterm{def:g10:lewis-and-shape:lone-pair}{pares no
enlazantes}, angular.
\end{solution}

\begin{solution}{exo:g10:lewis-and-shape:9}
En el \ce{CO2}, el carbono tiene dos grupos (dos \omterm{def:g10:lewis-and-shape:multiple-bond}{enlaces dobles}) y
ningún \omterm{def:g10:lewis-and-shape:lone-pair}{par no enlazante}: apuntan en sentidos opuestos. En el \ce{H2O},
el oxígeno tiene cuatro grupos (dos enlaces y dos \omterm{def:g10:lewis-and-shape:lone-pair}{pares no enlazantes})
en disposición tetraédrica: los dos enlaces forman un ángulo.
\end{solution}

\begin{solution}{exo:g10:lewis-and-shape:10}
\chemfig{C(-[3]H)(-[5]H)=C(-[1]H)-[7]H}: alrededor de cada carbono, tres
grupos (un \omterm{def:g10:lewis-and-shape:multiple-bond}{enlace doble} y dos sencillos): triangular plana, con ángulos
cercanos a \ang{120}.
\end{solution}

\begin{solution}{exo:g10:lewis-and-shape:11}
El carbono en el centro, con cuatro enlaces sencillos hacia \omterm{def:g7:atoms-and-molecules:atom}{átomos} de
cloro que tienen cada uno tres \omterm{def:g10:lewis-and-shape:lone-pair}{pares no enlazantes}: cuatro grupos
alrededor del carbono, tetraédrica.
\end{solution}

\begin{solution}{exo:g10:lewis-and-shape:12}
El etanol, \chemfig{H-C(-[2]H)(-[6]H)-C(-[2]H)(-[6]H)-\charge{90=\:,270=\:}{O}-H}, y el
metoximetano,
\chemfig{H-C(-[2]H)(-[6]H)-\charge{90=\:,270=\:}{O}-C(-[2]H)(-[6]H)-H}.
\end{solution}

\begin{solution}{exo:g10:lewis-and-shape:13}
Un \omterm{def:g10:lewis-and-shape:lone-pair}{par no enlazante} ocupa algo más de espacio que un \omterm{def:g10:lewis-and-shape:lone-pair}{par enlazante} y
acerca los enlaces entre sí. El amoniaco tiene un \omterm{def:g10:lewis-and-shape:lone-pair}{par no enlazante}, y el
agua, dos: el ángulo disminuye del metano al amoniaco y del amoniaco al
agua.
\end{solution}

\begin{solution}{exo:g10:lewis-and-shape:14}
\chemfig{H-C~\charge{0=\:}{N}}: dos grupos alrededor del carbono, lineal.
\chemfig{H-C(-[6]H)=\charge{45=\:,315=\:}{O}}: tres grupos alrededor del carbono,
triangular plana.
\end{solution}

\begin{solution}{exo:g10:lewis-and-shape:15}
$3 \times 6 = 18$ \omterm{def:g9:inside-the-atom:nucleus}{electrones}, 9 pares:
\chemfig{\charge{180=\:,270=\:}{O}=\charge{90=\:}{O}-\charge{0=\:,90=\:,270=\:}{O}}. El oxígeno central
tiene tres grupos (un \omterm{def:g10:lewis-and-shape:multiple-bond}{enlace doble}, uno sencillo y un \omterm{def:g10:lewis-and-shape:lone-pair}{par no
enlazante}): la \omterm{def:g7:atoms-and-molecules:molecule}{molécula} es angular.
\end{solution}

\begin{solution}{pb:g10:lewis-and-shape:1}
\textbf{1.} \chemfig{H-C(-[2]H)(-[6]H)-H}, \chemfig{H-\charge{90=\:}{N}(-[6]H)-H},
\chemfig{H-\charge{90=\:,270=\:}{O}-H}.

\textbf{2.} Cuatro grupos en cada una.

\textbf{3.} Ninguno en el metano, uno en el amoniaco y dos en el agua.

\textbf{4.} El carbono forma cuatro enlaces; el oxígeno forma dos (sus
otros dos grupos son \omterm{def:g10:lewis-and-shape:lone-pair}{pares no enlazantes}, que un juego de modelos no
representa con varillas).

\textbf{5.} \chemfig{H-C(-[2]H)(-[6]H)-C(-[2]H)(-[6]H)-H}: una varilla une
las dos bolas de carbono, y hacen falta seis bolas de hidrógeno.

\textbf{6.} Tetraédrica, piramidal, angular.

\textbf{7.} Los \omterm{def:g10:lewis-and-shape:lone-pair}{pares no enlazantes} ocupan más espacio que los \omterm{def:g10:lewis-and-shape:lone-pair}{pares
enlazantes} y aprietan los enlaces entre sí; el agua, con dos \omterm{def:g10:lewis-and-shape:lone-pair}{pares no
enlazantes}, tiene el ángulo más pequeño.

\textbf{8.} Los agujeros de una bola de carbono forman \ang{109.5}; el
ángulo real del agua es \ang{104.5}: una bola de oxígeno necesita sus
propios dos agujeros a \ang{104.5}.

\textbf{9.} \ang{180}: dos grupos alrededor del carbono apuntan en
sentidos opuestos, y la \omterm{def:g7:atoms-and-molecules:molecule}{molécula} es lineal.

\textbf{10.} Dos vértices que no están en la misma arista son vértices
opuestos de una cara: su distancia es la diagonal de la cara,
$\sqrt{a^2 + a^2} = a\sqrt{2}$.

\textbf{11.} El centro está en la mitad de una diagonal del cubo:
$a\sqrt{3}/2$.

\textbf{12.} H--H: $\sqrt{2} \approx \qty{1.41}{cm}$; C--H:
$\sqrt{3}/2 \approx \qty{0.87}{cm}$.

\textbf{13.} Cada hidrógeno está en un vértice del cubo y el carbono en
su centro: todos los vértices están a la misma distancia del centro.

\textbf{14.} El triángulo C\,H$_1$\,H$_2$ es isósceles (C\,H$_1$ =
C\,H$_2$); la recta que va de su vértice al punto medio de la base es
perpendicular a la base.

\textbf{15.} $\sin \widehat{\text{M\,C\,H}_1} = \dfrac{\text{M\,H}_1}{\text{C\,H}_1}
= \dfrac{a\sqrt{2}/2}{a\sqrt{3}/2} = \sqrt{2/3} \approx 0.816$.

\textbf{16.} Semiángulo $\approx \ang{54.7}$; ángulo H--C--H $\approx
2 \times 54.7 = \ang{109.5}$.

\textbf{17.} Es el ángulo de \ang{109.5} del tetraedro de la tabla.

\textbf{18.} Los agujeros hay que taladrarlos formando \ang{109.5} entre
sí.
\end{solution}
